\maketitle
\thispagestyle{firstpage}
\vspace{0.35em}

\begin{abstract}
Six Birds theory studies finite systems observed through deterministic coarse-grainings (``lenses'') using a small audited vocabulary of primitives.
This paper proves eight no-go theorems that constrain when common emergence narratives can be made honest in a finite, checkable setting.
For a finite Markov system and horizon $T$, define an \emph{arrow audit} as the Kullback--Leibler divergence between the observed path law and its time reversal.
We prove a deterministic data-processing inequality: observation cannot increase arrow.
As a consequence, any reversible system started in stationarity has zero observed arrow, even when the observation supports nontrivial protocol structure.
Further results give graph-theoretic obstructions for force-like antisymmetric drives (exactness on forests and zero cycle sums), a finite variational identity for best macro-kernels and closure deficit, and a bounded-interface saturation theorem that prevents unbounded laddering.
Together these theorems delimit what can and cannot be certified as emergent from audited primitives in finite models.
\end{abstract}

\vspace{0.2em}
\noindent\textbf{Keywords:} coarse-graining; finite Markov chains; data processing; exact graph forms; closure deficit; bounded interfaces.

\medskip

\noindent\textbf{Contribution summary.} This paper contributes an audited finite semantics for eight no-go fronts, a unified theorem framework linking arrow, graph exactness, closure, objecthood, and ladder obstructions, and a technical appendix structure that isolates routine derivations from the main proof flow.
